%-------------------------
% Cover Letter - Upstart Software Engineer, Borrower Experience
% Author: Ajay Venkatesh
%-------------------------

\documentclass[letterpaper,11pt]{article}

\usepackage[top=0.8in, bottom=0.8in, left=0.9in, right=0.9in]{geometry}
\usepackage[hidelinks]{hyperref}
\usepackage{lmodern}
\usepackage{parskip}
\usepackage{microtype}

\setlength{\parskip}{6pt}
\setlength{\parindent}{0pt}

\begin{document}

\begin{flushleft}
\textbf{\large Ajay Venkatesh} \\
Brooklyn, NY \\
+1 516 585 9013 \\
\href{mailto:av3855@nyu.edu}{av3855@nyu.edu}
\end{flushleft}

\vspace{10pt}

May 10, 2026

\vspace{6pt}

Hiring Manager \\
Upstart

\vspace{10pt}

Dear Hiring Manager,

My name is Ajay Venkatesh --- finishing my M.S. in Computer Engineering at NYU, with several years of production software experience: a few years at PwC in Bengaluru, a summer SDE internship at Amazon, and ongoing on-campus work at NYU's Novel AI Technologies. My focus has been on \textbf{full-stack product systems, distributed backends,} and the practical \textbf{AI layer} --- LLM-powered features, RAG pipelines, and agentic workflows that ship in real products.

The Upstart Borrower Experience posting hits a few of my strongest threads. At NYU's Novel AI Technologies I shipped a customer-facing insurance analytics platform end-to-end --- \textbf{React, Node.js, REST APIs}, \textbf{Docker} on \textbf{EC2}, with Stripe billing, role-based access, and row-level security across paying clients. The shape of that work --- borrower-style self-service, financial workflows, and reducing friction in the user journey --- maps directly onto what the posting describes. I also led a broker-portal integration with insurer platforms using \textbf{document/media processing} and \textbf{WebSocket-based real-time communication} --- the same kind of "translate operational workflows into intuitive product experiences" the role calls out.

On AI: at Amazon I deployed a \textbf{Retrieval-Augmented Generation (RAG)} pipeline with an \textbf{agentic workflow} that autonomously retrieves operational logs and generates context-aware summaries and remediation steps. I implemented an \textbf{MCP server} for tool-augmented retrieval and built a high-throughput, event-driven triage system on \textbf{ECS Fargate, SQS, SNS} that materially reduced response latency. At Campus Mesh I built \textbf{CampusIQ}, an LLM-powered AI assistant with \textbf{prompt engineering} and \textbf{RAG retrieval} over structured data, with embedding-based semantic search to improve contextual accuracy. Building LLM-powered features in user-facing surfaces, with safety and explainability in mind, is exactly where I want to spend the next chapter.

Stack-wise: \textbf{Python} and \textbf{TypeScript / JavaScript} for backend and product surfaces; \textbf{React} for web; cloud microservices on \textbf{AWS} with \textbf{Docker} and CI/CD; \textbf{SQL} (PostgreSQL, MySQL, DynamoDB) for the data layer. I'm comfortable spanning the stack, partnering with product / design / analytics, and reasoning about reliability and safety in customer-facing systems.

What pulls me toward Upstart specifically is the mission --- reducing the cost and complexity of borrowing through AI --- and the fact that the work is measured by real outcomes for real borrowers, not internal vanity metrics. That's the kind of feedback loop I want to be inside of.

I'd welcome the chance to discuss how I could contribute to the Borrower Experience team. Thanks for considering my application.

\vspace{12pt}

Sincerely, \\
Ajay Venkatesh

\end{document}
